\section{Related work}
\textbf{T2I safety benchmarks.}
Early benchmarks like I2P~\citep{schramowski2023safe} (4,703 prompts, 7 categories) and Unsafe Diffusion~\citep{qu2023unsafe} (434 prompts) sourced real user prompts from Lexica to study NSFW generation, while SneakyPrompt~\citep{yang2023sneakyprompt} introduced 200 adversarially crafted prompts.
Recent efforts~\citep{dai2024safesora,miao2024t2vsafetybench,li2025t2isafety,zhang2025t2iriskyprompt} expanded coverage to more risk dimensions.
Most related to our work, T2ISafety~\citep{li2025t2isafety} established a broad framework assessing fairness, toxicity, and privacy across $\sim$70K prompts, and T2I-RiskyPrompt~\citep{zhang2025t2iriskyprompt} introduced a hierarchical taxonomy of 6 primary and 14 fine-grained categories with reasoning-based annotation.
Complementary image-safety benchmarks such as UnsafeBench~\citep{yiting2025unsafebench} evaluate classifiers over real-world and AI-generated images but retain binary labels, leaving severity-graded generation risk unmeasured.
Both, however, still assign each prompt a single, independently sampled binary verdict: neither pairs multiple severity instances of the same scenario under matched generation conditions, so neither can compute how a defense's suppression at one severity relates to its suppression at another for a fixed scenario.
This distinction is not cosmetic: it is precisely the quantity Insight~3 requires.
\dataname's paired ladders hold the scenario fixed (same T2I model and seed across all four rungs); a flat prompt set cannot, making the relative-suppression asymmetry structurally invisible to it.
Our evaluation matrix (models, defenses, filters, attacks) otherwise follows the protocol of contemporary work for baseline consistency.

\noindent
\textbf{Defenses and filters.}
Internal defenses modify model weights~\citep{gandikota2023concept,gandikota2024unified} or retrain safety objectives~\citep{liu2024safetydpo,lu2024mace,chen2025trce}; external ones apply keyword lists~\citep{george}, text classifiers~\citep{li_NSFW_text_classifier,Detoxify,khader2024diffguard}, image-level detectors~\citep{laionAI,Chhabra,schramowski2023safe,Giphy}.
They all target L3-equivalent explicit content under a binary safe/unsafe protocol.
\dataname provides the first graded diagnostic across all four severity rungs, revealing that defenses effective at L3 remain vulnerable to low-to-mid-risk escalation.
